\subsection{SUBRINGS}

DEFINITION 3. Let A be a ring. A subring of A is any subset B of A which is a subgroup of A under addition, which is stable under multiplication and which contains the unit of A.

The above conditions may be written as follows

\[
0 \in \mathbf {B}, \quad \mathbf {B} + \mathbf {B} \subset \mathbf {B}, \quad - \mathbf {B} \subset \mathbf {B}, \quad \mathbf {B}. \mathbf {B} \subset \mathbf {B}, \quad 1 \in \mathbf {B}.
\]

If B is a subring of A, it is given the addition and multiplication induced by those on A, which make it into a ring. The canonical injection of B into A is a ring homomorphism.

Examples. (1) Every subgroup of the additive group \(\mathbf{Z}\) which contains 1 is equal to \(\mathbf{Z}\) . Thus \(\mathbf{Z}\) is the only subring of \(\mathbf{Z}\) .

\begin{enumerate}
\def\labelenumi{(\arabic{enumi})}
\setcounter{enumi}{1}
\item
  Let A be a ring and \((A_{t})_{t\in I}\) a family of subrings of A; it is immediate that \(\bigcap_{i\in I}A_{i}\) is a subring of A. In particular, the intersection of the subrings of A containing a subset X of A is a subring called the subring of A generated by X.
\item
  Let X be a subset of a ring A. The centralizer of X in A is a subring of A. In particular, the centre of A is a subring of A.
\item
  Let G be a commutative group with operators; let \(\Omega\) denote the set of operators and \(\alpha \mapsto f_{\alpha}\) the action of \(\Omega\) on G. Let E be the endomorphism ring of the group without operators G and F the set of endomorphisms of the group with operators G. By definition, F consists of the endomorphisms \(\phi\) of G such that \(\phi \cdot f_{\alpha} = f_{\alpha} \cdot \phi\) for all \(\alpha \in \Omega\) . Therefore F is a subring of the ring E. F is called the endomorphism ring of the group with operators G (cf.~II, § 1, no. 2). Let \(F_{1}\) be the subring of E generated by the \(f_{\alpha}\) . Then F is the centralizer of \(F_{1}\) in E.
\end{enumerate}

